\documentclass[11pt]{article}

\usepackage{amsmath,amssymb,amsthm}
\usepackage[margin=1.1in]{geometry}
\usepackage{booktabs}
\usepackage{algorithm}
\usepackage{algpseudocode}

\newtheorem{theorem}{Theorem}[section]
\newtheorem{proposition}[theorem]{Proposition}
\newtheorem{lemma}[theorem]{Lemma}
\newtheorem{corollary}[theorem]{Corollary}
\newtheorem{conjecture}[theorem]{Conjecture}
\theoremstyle{definition}
\newtheorem{definition}[theorem]{Definition}
\newtheorem{remark}[theorem]{Remark}
\newtheorem{measured}[theorem]{Measured Result}

\newcommand{\C}{\mathbb{C}}
\newcommand{\R}{\mathbb{R}}
\newcommand{\N}{\mathbb{N}}
\newcommand{\cA}{\mathcal{A}}
\newcommand{\cM}{\mathcal{M}}
\newcommand{\cL}{\mathcal{L}}
\newcommand{\ip}[2]{\langle #1,\,#2\rangle}
\newcommand{\norm}[1]{\lVert #1\rVert}
\newcommand{\abs}[1]{\lvert #1\rvert}
\newcommand{\diag}{\operatorname{diag}}
\newcommand{\supp}{\operatorname{supp}}

\title{Super-resolution by intersection of aperture-magnitude varieties:\\
a formal account of the two-lattice / aperture-ladder method}
\author{BFFT project working notes}
\date{2026-07-11 \; (v1)}

\begin{document}
\maketitle

\begin{abstract}
We formalize the super-resolution method deployed in the BFFT viewer and its
image-fusion generalization.  One latent signal is constrained by the
\emph{magnitudes} of several linear analyses (``apertures'': STFT windows of
different lengths, blur channels, focus slices).  The estimate is obtained by
alternating closed-form projections onto the corresponding nonconvex
magnitude varieties from a data-driven seed.  We prove that the synthesis
step used in practice is the Moore--Penrose pseudo-inverse of the analysis
frame, that the true signal is an exact fixed point of every projector, and
that the sensitivity of each family to a component-phase rotation (the
quantity that decides what a family can super-resolve) admits a sharp
closed-form bound: the \emph{min-overlap} of the component spectrograms.
Empirically this bound is tight up to a universal factor
$1/\sqrt2$ (log--log correlation $0.9995$, slope $0.997$), turning aperture
design into a computable geometry problem.  We define coverage, dead zones,
and aperture ladders; state the measured critical-ratio law for geometric
ladders ($r^{*}\in(4,8]$); relate local convergence rates to the overlap
functional; formalize relaxed (reliability-weighted) projections for
mutually inconsistent families; and delimit the method against the
optimal-transport barycenter formulation of Valdivia, Cazelles and
F\'evotte~\cite{valdivia2026}, from which this work departed.  Statements
are tagged as \emph{proved}, \emph{measured}, or \emph{conjectured}.
\end{abstract}

\section{Introduction}

Time--frequency super-resolution by fusing analyses of different window
lengths was recently formalized in~\cite{valdivia2026} as an unbalanced
optimal-transport (UOT) barycenter problem over spectrogram measures.  That
formulation is convex in the transport plans and is solved by a block
majorization--minimization scheme; its output is a measure on a
time--frequency grid.

The present method replaces the barycentric \emph{average of
representations} by the \emph{intersection of measurement constraints on a
latent signal}.  The observed spectrogram magnitudes are not objects to be
transported; they are equations that the latent signal must satisfy.  All
super-resolving information then travels through the phase variables that
magnitudes do not record, and the algorithmic content reduces to alternating
closed-form projections.  Three structural consequences, each proved or
measured below, distinguish the two formulations:
(i) the true signal is an \emph{exact fixed point} of our iteration
(Proposition~\ref{prop:fixed}), whereas a barycenter is generically not the
spectrogram of any signal (Remark~\ref{rem:consistency});
(ii) what can and cannot be super-resolved is decided by a computable
overlap functional of the apertures (Theorem~\ref{thm:overlap},
Measured Result~\ref{meas:tight}), not by a transport cost;
(iii) the per-tile cost is a fixed, small number of FFT passes
(Proposition~\ref{prop:complexity}), compatible with streaming real-time
operation.

\section{Notation and preliminaries}

Signals are vectors $u\in\C^{L}$.  For $N\in\N$ dividing conventions we use
the \emph{unitary} discrete Fourier transform $F_{N}\in\C^{N\times N}$,
$(F_{N})_{mn}=N^{-1/2}e^{-2\pi i mn/N}$, so that $F_{N}^{*}F_{N}=I$.  (The
implementation uses the unnormalized transform with compensating factors;
all statements are invariant to this choice.)

Fix a window $w\in\R^{N}$ (symmetric Hann throughout), a hop $H\in\N$, and
the frame offsets
\[
  \tau_{k}=kH,\qquad k=0,\dots,K-1,\qquad K=\lfloor (L-N)/H\rfloor+1 .
\]
Let $S_{k}:\C^{L}\to\C^{N}$ be the selection $(S_{k}u)_{n}=u_{\tau_{k}+n}$.

\begin{definition}[Aperture / analysis operator]\label{def:aperture}
The \emph{aperture} with parameters $(N,H,w)$ is the linear map
\[
  \cA:\C^{L}\to\C^{K\times N},\qquad
  (\cA u)_{k}=F_{N}\bigl(w\odot S_{k}u\bigr).
\]
A finite collection $\{\cA_{i}\}_{i=1}^{P}$ of apertures (possibly of
different lengths $N_{i}$, hops $H_{i}$, or acting in other linearly
reachable domains, cf.\ \S\ref{sec:domains}) is called a \emph{family set};
when the $N_{i}$ form a scale sequence we call it an \emph{aperture ladder}.
\end{definition}

\begin{definition}[Weight function and COLA support]
$d\in\R^{L}$, $d(n)=\sum_{k} w^{2}(n-\tau_{k})$ (indices where defined).
The \emph{support} of the aperture is $\supp(\cA)=\{n: d(n)>0\}$.  The
aperture is \emph{complete} if $\supp(\cA)=\{0,\dots,L-1\}$.
\end{definition}

Throughout, magnitudes, products and quotients of coefficient arrays are
taken entrywise, and $\norm{\cdot}$ is the Euclidean ($\ell^{2}$ /
Frobenius) norm.

\section{Measurement model and constraint varieties}

Let $z\in\C^{L}$ be the latent signal and let
\[
  m_{i}\;=\;\abs{\cA_{i}z}\in\R_{\ge0}^{K_{i}\times N_{i}},
  \qquad i=1,\dots,P,
\]
be the observed magnitude arrays.  (In the audio deployment the $m_i$ are
computed from the recorded IQ tile; in image fusion each capture supplies
the $m_i$ of its own forward operator.)  Phases of $\cA_i z$ are
\emph{not} used by the iteration; the observation phase enters exactly once,
to fix the global gauge (\S\ref{sec:gauge}).

\begin{definition}[Magnitude variety]
$\cM_{i}=\{u\in\C^{L}: \abs{\cA_{i}u}=m_{i}\}$.  The \emph{feasible set} of
the family set is $\cM=\bigcap_{i=1}^{P}\cM_{i}$.
\end{definition}

Each single-coefficient constraint $\abs{y_c}=m_c$ is a circle in $\C$;
$\cM_i$ is the preimage of a torus under a linear map, a real-analytic
variety that is nonconvex whenever some $m_c>0$.  By construction
$z\in\cM$: the feasible set is never empty, and super-resolution is the
statement that $\cM$ is \emph{small} around $z$ modulo the gauges of
\S\ref{sec:gauge}.

\begin{remark}[No convex signal-space model]\label{rem:convex}
If $u\in\cM$ then $e^{i\theta}u\in\cM$ for all $\theta$; any convex set
containing this orbit contains $0\notin\cM$ (whenever some $m_c>0$).  Hence
no convex subset of $\C^{L}$ equals $\cM$; convexification requires the
lifted matrix variable of PhaseLift~\cite{candes2013}, at dimension
$L^{2}$.  The method below instead exploits that $\cM$ is
\emph{locally} benign near $z$ (\S\ref{sec:convergence}).
\end{remark}

\section{Projectors}\label{sec:projectors}

\begin{definition}[Radial map and family projector]\label{def:projector}
For a target array $m$, define the radial map
$R(y)=m\odot y/\abs{y}$ (entrywise, with $R(y)_c=0$ when $m_c=0$ and
$R(y)_c$ left equal to $y_c$ scaled to $m_c$ arbitrarily on the null set
$y_{c}=0$).  The \emph{family projector} of aperture $i$ is
\[
  P_{i}(u)\;=\;\cA_{i}^{+}\,R_{i}\bigl(\cA_{i}u\bigr),
\]
where $\cA_{i}^{+}$ denotes the Moore--Penrose pseudo-inverse, and samples
outside $\supp(\cA_i)$ carry the input value ($P_i(u)_n=u_n$ there).
\end{definition}

\begin{lemma}[The OLA synthesis is the pseudo-inverse]\label{lem:pinv}
For a complete aperture, $\cA^{*}\cA=\diag(d)$ with $d>0$, and
\[
  \cA^{+}Y=\diag(d)^{-1}\cA^{*}Y
  =\Bigl(\textstyle\sum_{k}S_{k}^{T}\bigl(w\odot F_{N}^{*}Y_{k}\bigr)\Bigr)
  \big/\,d ,
\]
i.e.\ inverse DFT of each frame, windowing, overlap--add, and division by
the accumulated squared window.  This is exactly the synthesis implemented
in \texttt{two\_lattice.MagnitudeFamily.project} and
\texttt{MagnitudeProjectorC::project}.
\end{lemma}

\begin{proof}
$\cA^{*}Y=\sum_{k}S_{k}^{T}(w\odot F_{N}^{*}Y_{k})$ by unitarity of
$F_{N}$ and the definition of adjoints of $S_k$ and of entrywise
multiplication by the real window.  Hence
$\cA^{*}\cA u=\sum_{k}S_{k}^{T}(w^{2}\odot S_{k}u)=d\odot u$.  For an
injective operator ($d>0$), $\cA^{+}=(\cA^{*}\cA)^{-1}\cA^{*}$.
\end{proof}

\begin{remark}
$P_i$ is the classical Griffin--Lim composition~\cite{griffinlim1984}:
$R_i$ is the exact metric projection of the coefficient array onto the
torus $\{\abs{y}=m_i\}$ (nearest point on each circle), and
$\cA_i\cA_i^{+}$ is the orthogonal projection onto the \emph{consistent}
subspace $\operatorname{range}\cA_i$.  $P_i$ is a surrogate for---not equal
to---the metric projection onto $\cM_i$; for the single-family iteration
$u\mapsto P_i(u)$ the coefficient-domain residual
$\norm{\,\abs{\cA_i u}-m_i\,}$ is nonincreasing~\cite{griffinlim1984}.  No
such monotonicity is claimed for alternation between families.
\end{remark}

\begin{proposition}[Exact fixed point]\label{prop:fixed}
If $m_{i}=\abs{\cA_{i}z}$ for all $i$, then $P_{i}(z)=z$ for every complete
aperture, and $z$ is a fixed point of any composition, relaxation, or
momentum extrapolation of the $P_{i}$ that fixes common fixed points.
\end{proposition}

\begin{proof}
$R_{i}(\cA_{i}z)=\abs{\cA_i z}\odot \cA_i z/\abs{\cA_i z}=\cA_{i}z$, and by
Lemma~\ref{lem:pinv} $\cA_{i}^{+}\cA_{i}=\diag(d)^{-1}\diag(d)=I$.
\end{proof}

Proposition~\ref{prop:fixed} was verified numerically to machine precision
and separates this formulation from any averaging scheme: the truth is not
merely near-optimal, it is \emph{stationary}.  Consequently the deployed
one-iteration algorithm (\S\ref{sec:algorithm}) is a local polisher of its
seed toward $\cM$, and all questions of super-resolution reduce to the local
geometry of $\cM$ at $z$.

\section{Gauge geometry and the overlap theorem}\label{sec:gauge}

Write $z=a+b$ for a decomposition into two components (e.g.\ two tones, two
packets, two image structures), and consider the \emph{relative-phase
orbit} $z_{\varphi}=a+e^{i\varphi}b$.  Whether a family can distinguish
$z_\varphi$ from $z$ is precisely whether it can super-resolve the pair,
since the relative phase is what encodes beat timing / sub-cell structure.

\begin{definition}[Pin strength]\label{def:pin}
The pin strength of aperture $\cA$ on the pair $(a,b)$ at angle $\varphi$
is
\[
  E_{\varphi}(a,b;\cA)\;=\;
  \frac{\bigl\lVert\,\abs{\cA z_{\varphi}}-\abs{\cA z}\,\bigr\rVert}
       {\norm{\cA z}} .
\]
A direction with $E\equiv0$ for all families is a \emph{gauge} of the
family set; the pair is \emph{pinned} by the set if some family has
$E_{\varphi}>0$.
\end{definition}

\begin{theorem}[Min-overlap bound]\label{thm:overlap}
For any aperture $\cA$ and any $\varphi$,
\[
  E_{\varphi}(a,b;\cA)\;\le\;
  \abs{1-e^{i\varphi}}\;
  \frac{\Bigl(\sum_{c}\min\bigl(\abs{(\cA a)_{c}},\abs{(\cA b)_{c}}\bigr)^{2}
        \Bigr)^{1/2}}
       {\norm{\cA z}},
\]
where $c$ ranges over all time--frequency cells and
$\abs{1-e^{i\varphi}}=2\abs{\sin(\varphi/2)}$.
In particular, if the cell supports of $\cA a$ and $\cA b$ are disjoint
(the min vanishes in every cell), the relative-phase orbit lies entirely in
$\cM_{\cA}$: the pair is an exact gauge of that family.
\end{theorem}

\begin{proof}
Per cell, with $\alpha=(\cA a)_{c}$ and $\beta=(\cA b)_{c}$, the reverse
triangle inequality gives
$\bigl|\,\abs{\alpha+e^{i\varphi}\beta}-\abs{\alpha+\beta}\,\bigr|
 \le\abs{e^{i\varphi}\beta-\beta}=\abs{1-e^{i\varphi}}\,\abs{\beta}$.
Factoring the unimodular $e^{i\varphi}$ from the first term,
$\abs{\alpha+e^{i\varphi}\beta}=\abs{e^{-i\varphi}\alpha+\beta}$, the same
inequality gives the bound with $\abs{\alpha}$ in place of $\abs{\beta}$.
Taking the smaller of the two and summing squares over cells yields the
claim; the disjoint-support case makes every summand zero, for every
$\varphi$, hence $z_{\varphi}\in\cM_{\cA}$.
\end{proof}

\begin{measured}[Tightness up to a universal factor]\label{meas:tight}
Over $57$ packet-pair configurations spanning three decades of pin strength
(separations $\Delta t\in[60,2200]$ samples, $\Delta f\in[0,12]$ bins,
windows $N\in\{256,\dots,2048\}$; \texttt{experiments/ladder\_theorem.py}),
the measured $E_{\pi/2}$ against the bound of Theorem~\ref{thm:overlap}
satisfies: log--log Pearson correlation $0.9995$, regression slope $0.997$,
ratio median $0.685$ with interquartile range $[0.673,0.707]$.  The ratio
is consistent with the value $1/\sqrt2$ obtained by averaging the radial
component of the rotated interference over alignment angles.  Hence, to
within a few percent,
\[
  E_{\pi/2}(a,b;\cA)\;\approx\;
  \Bigl(\sum_{c}\min\bigl(\abs{\cA a}_{c},\abs{\cA b}_{c}\bigr)^{2}\Bigr)^{1/2}
  \Big/\;\norm{\cA z} .
\]
The right-hand side is computable from single-component spectrograms alone:
aperture design requires no probing.
\end{measured}

\begin{corollary}[Benign gauges of magnitude displays]
Components that are cell-disjoint under \emph{every} family of the set
(e.g.\ an isolated transient far from all other content) retain a free
absolute phase.  This gauge is invisible to any magnitude readout and is
therefore harmless for display purposes; it matters only for
waveform-domain post-processing, where an additional constraint must pin
it.  (Measured residuals for an isolated click: $0.008$ and $0.002$ for the
deployed pair, i.e.\ numerically at the probe floor.)
\end{corollary}

\begin{measured}[Complementarity of the deployed pair]\label{meas:pair}
For a tone pair separated by $1.4$ frequency bins of the long window
($N_B=1024$), the measured pin strengths are $0.264$ (long) and $0.522$
(short, $N_S=256$): the \emph{short} family pins the pair's relative phase
twice as hard, because both tones occupy one short-window cell and the
interference (beat) term appears in its magnitudes, while the long window
separates them into different cells.  Conversely only the long family
resolves the pair's existence.  This complementarity of constraints on the
same latent degrees of freedom is the super-resolution mechanism; no
representation-space blending occurs anywhere.
\end{measured}

\section{Coverage, dead zones, ladders}\label{sec:coverage}

\begin{definition}[Coverage of a family set]
For a content class parameterized by pair separations
$(\Delta t,\Delta f)$ at component width $w$, the coverage of the set
$\Lambda=\{\cA_{i}\}$ is
$\sigma_{\Lambda}(\Delta t,\Delta f)=\max_{i}E_{\pi/2}(a,b;\cA_{i})$
evaluated on canonical packet pairs.  The \emph{pinnable region} is where
$\sup_{N}\sigma_{\{\cA_N\}}>\varepsilon$ over all admissible windows; the
\emph{dead zone} of $\Lambda$ is the pinnable region minus
$\{\sigma_{\Lambda}>\varepsilon\}$.
\end{definition}

\begin{proposition}[Visibility law; support geometry]\label{prop:visibility}
For Hann packets of width $w$ at separation $(\Delta t,\Delta f)$, an
aperture of length $N$ has $\sigma_N>0$ essentially iff some analysis cell
contains energy of both components, i.e.
\[
  \Delta t + w \;\lesssim\; N
  \qquad\text{and}\qquad
  \Delta f \;\lesssim\; \tfrac{c_1}{w}+\tfrac{c_2}{N},
\]
with constants of order one determined by the window's essential supports.
Consequently a pair with
$(\Delta f - c_1/w)\,(\Delta t + w)\gtrsim 1$ is invisible to \emph{every}
aperture (the uncertainty boundary), and coverage maps of single rungs are
rectangles in $(\Delta t,\Delta f)$ with strength decaying toward the cell
edges.  \textup{(Statement at the level of essential supports; constants
measured in \texttt{experiments/aperture\_ladder\_coverage.py}, where the
heatmaps realize exactly this geometry.)}
\end{proposition}

\begin{measured}[Deployed-pair coverage and interior redundancy]
On a $(\Delta t,\Delta f)$ grid with $w=200$, $L=8192$: the deployed ladder
$\{1024,256\}$ covers $53.2\%$ of the pinnable set; adding a $2048$ rung
covers $100\%$; adding an interior $512$ rung instead adds exactly nothing.
The deployed ladder's entire dead zone is the band
$\Delta t\in(800,1500)$.
\end{measured}

\begin{measured}[Critical ratio of geometric ladders]\label{meas:rstar}
Along the $\Delta f=0$ ridge, relative to the dense ($r=2$) strength
envelope, geometric ladders with ratio $r=4$ attain the envelope
\emph{exactly} (hole ratio $1.000$) and ladders with $r=8$ possess points
of ratio $0.000$ (total holes), for packet widths $w\in\{100,200,400\}$ and
span $[256,4096]$.  Hence the critical ratio satisfies $r^{*}\in(4,8]$ for
this content class, and the deployed geometry $N_B/N_S=4$ is critical-safe.
Covering a scale span $[N_{\min},N_{\max}]$ requires
$\lceil\log_{4}(N_{\max}/N_{\min})\rceil+1$ rungs.
\end{measured}

\begin{conjecture}[Ladder theorem]\label{conj:ladder}
For components of essential width $w$ and window family
$\{\cA_{N}\}$, define the single-rung strength profile
$s_{N}(\Delta t)=\sigma_{\{\cA_N\}}(\Delta t,0)$.  There is a critical
ratio $r^{*}(w,\varepsilon)$, determined by the crossing points of
neighboring profiles via the min-overlap functional of
Theorem~\ref{thm:overlap}, such that a geometric ladder of ratio
$r\le r^{*}$ attains the dense envelope up to factor $1-\varepsilon$
throughout its scale span, while some $r>r^{*}$ ladder has interior zeros.
The measured value $r^{*}\in(4,8]$ (Measured Result~\ref{meas:rstar}) and
the closed-form profile of Measured Result~\ref{meas:tight} reduce the
proof to elementary support estimates for Hann windows.
\end{conjecture}

\section{Convergence}\label{sec:convergence}

The iteration is (heavy-ball accelerated) alternating projection between
the $\cM_i$.  Local linear convergence of alternating projections for
nonconvex sets holds under transversality of the intersection, with rate
governed by the angle between tangent spaces; see
\cite{lewisluke2009,drusvyatskiy2015}.  We do not restate those theorems;
we record the empirical identification that makes them quantitative here.

\begin{measured}[Strength--rate law]\label{meas:rate}
Restart a pair on its relative-phase orbit ($\varphi_{0}=43^{\circ}$,
$\Delta t=1600$) and iterate.  With pin strength $0.000$ (pair not covered)
the error is retained indefinitely ($43.1^{\circ}$ after $300$ iterations);
with strength $0.016$ (covered at a cell edge) the error decays to
$13.1^{\circ}$ after $300$; with strength $0.286$ (covered deep in a cell,
$N=4096$) it decays to $0.8^{\circ}$ by iteration $20$.  The decay rate
scales approximately as the square of the pin strength, consistent with the
angle interpretation.  \emph{Design rule: place expected content deep in
some rung's cell (rung length ${\approx}2$--$2.5\times$ the content
separation); cell-edge coverage is coverage in principle only.}
\end{measured}

\begin{measured}[Seed and basin]\label{meas:seed}
From random initialization the alternating iteration stagnates (residual
ratio ${\approx}0.996$--$0.998$ per step), the known Griffin--Lim regime.
From a phase-degraded seed of fidelity $0.81$ it reaches $0.93$ in one
iteration and $0.99$ in $24$.  The deployed seed (magnitude-only PGHI
phase propagation~\cite{prusa2017} through the paused-DIP state, one
gradient step, consensus) lands inside this basin on recorded IQ tiles;
one unified iteration then suffices.  All global-search difficulty of the
nonconvex problem is thereby carried by the seed, not by the projections.
\end{measured}

\section{Relaxed projections for inconsistent families}\label{sec:relaxed}

When the families are measured by \emph{different} instruments
(multi-focus or multi-blur captures), the arrays $m_i$ are generally not
the magnitudes of any common latent: $\cM=\emptyset$.  The deployed
generalization weights the radial correction per coefficient.

\begin{definition}[Relaxed family projector]\label{def:relaxed}
For weights $\omega_i\in[0,1]^{K_i\times N_i}$,
\[
  P_{i}^{\omega}(u)
  \;=\;\cA_{i}^{+}\Bigl((1-\omega_{i})\odot\cA_{i}u
        +\omega_{i}\odot R_{i}(\cA_{i}u)\Bigr)
  \;=\;u+\cA_{i}^{+}\Bigl(\omega_{i}\odot
        \bigl(R_{i}(\cA_{i}u)-\cA_{i}u\bigr)\Bigr),
\]
the second form by Lemma~\ref{lem:pinv}.  $\omega\equiv1$ recovers
Definition~\ref{def:projector}; $\omega\equiv0$ ignores the family.
\end{definition}

\begin{proposition}[Stationarity of the averaged relaxed scheme]
A point $u$ is stationary for the simultaneous (averaged) relaxed
iteration $u\mapsto u+\tfrac1P\sum_{i}\cA_{i}^{+}(\omega_{i}\odot
(R_{i}-I)\cA_{i}u)$ iff
\[
  \sum_{i}\cA_{i}^{+}\Bigl(\omega_{i}\odot
   \bigl(m_{i}-\abs{\cA_{i}u}\bigr)\odot
   \tfrac{\cA_{i}u}{\abs{\cA_{i}u}}\Bigr)=0 ,
\]
i.e.\ the reliability-weighted magnitude residuals cancel after synthesis:
fixed points are weighted compromises exactly where families disagree.
\end{proposition}

\begin{proof}
Immediate from the second form in Definition~\ref{def:relaxed} and
$R_i(y)-y=(m_i-\abs{y})\odot y/\abs{y}$.
\end{proof}

\begin{remark}[Relation to unbalanced OT]
The weights $\omega_i$ play the role of the marginal-relaxation parameters
$\eta$ in the UOT formulation of~\cite{valdivia2026}: both interpolate
between enforcing a family exactly and ignoring it, and both make the
mutually-inconsistent case well-posed.  The difference is the ambient
space: UOT relaxes marginals of a transport plan between representations;
Definition~\ref{def:relaxed} relaxes constraint enforcement on the latent
signal.
\end{remark}

\begin{measured}[Reliability weights that work, and two that do not]
With $\omega$ = per-patch softmax of high-band energy across captures
(defocus is a local low-pass, so relative high-band energy is measured
reliability), the relaxed scheme fuses the MFFW3/4 focus stack without
ringing or ghosting and beats a hard arg-max baseline
(Laplacian variance $0.0101$ vs $0.0082$; best source $0.0040$).
Two measured cautions: (i) an \emph{absolute} gate modeled on noise floors
mis-weights coefficients that are confidently \emph{biased} (blurred), not
noisy; (ii) when \emph{no} capture is locally reliable the residual is
information-limited: Douglas--Rachford, RAAR, and gating all coincide with
plain alternation to $\pm0.07$ dB, and only an added capture (a third blur
angle) moved the dead-band error ($1.041\to0.813$; PSNR $26.30\to26.41$).
Conversely, in an optimization-limited instance (global Fourier-modulus
fusion, \S\ref{sec:domains}), Douglas--Rachford / RAAR
splitting~\cite{luke2005} gains $+2.4$ dB over alternation.  The DR/AP gap
is thus a practical \emph{discriminator} between the two failure classes.
\end{measured}

\section{Domain freedom}\label{sec:domains}

Definition~\ref{def:aperture} nowhere requires the aperture to act in the
time domain: any injective-on-support linear $\cA_i$ with a fast adjoint
qualifies, e.g.\ two-dimensional patch transforms (image fusion), or the
global Fourier transform, for which the radial map is the \emph{exact}
metric projection since $F$ is unitary
($P(u)=F^{*}(m\odot Fu/\abs{Fu})$).  Measured consequences: a
capture misregistered by an unknown circular shift can be fused through
its shift-invariant Fourier modulus (naive $12.9$ dB $\to$ DR fusion
$24.6$ dB); a real (non-periodic) shift breaks modulus invariance at the
aperture boundary; \emph{patched} Fourier moduli are not shift-invariant at
all (targets constrain wrong locations).  Full-support natural images lack
the compact-support constraint of classical phase
retrieval~\cite{fienup1982}, so the global-modulus family is weakly
identifying on its own; its practical role is diagnosis (phase
correlation), with spatial families carrying the fusion.

\section{The deployed algorithm}\label{sec:algorithm}

\begin{algorithm}[h]
\caption{Unified ladder solve (one tile; \texttt{iqw\_dip\_unified\_ladder})}
\begin{algorithmic}[1]
\State \textbf{Input:} tile $z\in\C^{L}$; ladder $\{(N_{i},H_{i})\}_{i=1}^{P}$
       ordered largest first; iterations $T$ (deployed: $1$);
       momentum $\beta=0.88$; terminal relaxation $\rho=0.75$
\State $m_{i}\gets\abs{\cA_{i}z}$ for all $i$
       \Comment{measured targets; phases discarded}
\State $u_{0}\gets$ DIP/PGHI seed from $m$ (magnitude-only phase
       propagation, one gradient step, consensus) \Comment{basin entry}
\For{$t=1,\dots,T$}
  \State $v\gets u_{t-1}+\beta\,(u_{t-1}-u_{t-2})$ \Comment{inert at $t{=}1$}
  \For{$i=1,\dots,P$} $v\gets P_{i}(v)$ \Comment{Def.~\ref{def:projector}}
  \EndFor
  \State $u_{t}\gets v$
\EndFor
\State $u\gets u_{T}+\rho\,(P_{1}(u_{T})-u_{T})$
       \Comment{empirical early-stop finish; see note \S1 audit}
\State $u\gets u\cdot\ip{u}{z}/\abs{\ip{u}{z}}$
       \Comment{global gauge fixed once}
\State \textbf{Output:} latent $u$; display = reassigned STFT of $u$
\end{algorithmic}
\end{algorithm}

\begin{proposition}[Complexity]\label{prop:complexity}
One pass of the projector cascade costs
$\sum_{i=1}^{P}\Theta\bigl((L/H_{i})\,N_{i}\log N_{i}\bigr)$ flops; with
the deployed hops $H_{i}=N_{i}/8$ this is $\Theta(P\,L\log N_{\max})$,
independent of the rung lengths individually.  Long rungs are the cheap
direction: at $L=8192$, adding an $N=4096$ rung ($9$ frames) to the
$N_B=1024$ pair changed the measured tile time from $15.8$ to $16.3$ ms.
For comparison, the block-MM UOT barycenter of~\cite{valdivia2026}
optimizes over transport plans with $10^{4}$--$10^{7}$ retained entries for
$57$--$945$ iterations (reported $0.43$--$9.4$ s on second-scale,
kilohertz-rate material).
\end{proposition}

\begin{remark}[Barycenters are generically inconsistent]\label{rem:consistency}
The set of \emph{consistent} magnitude arrays
$\abs{\cA(\C^{L})}\subset\R_{\ge0}^{K\times N}$ has real dimension at most
$2L-1$, while redundant apertures have $KN>2L$ coefficients; consistent
arrays therefore form a measure-zero subset, characterized by nontrivial
inter-frame relations.  A barycenter of two consistent arrays computed in
measure space has no reason to satisfy those relations, and generically
does not: it is a picture, not the analysis of any signal.  The intersection
formulation returns a signal by construction, and every displayed quantity
is a readout of it.
\end{remark}

\section{Summary of epistemic status}

\begin{center}
\begin{tabular}{@{}lll@{}}
\toprule
Statement & Status & Where\\
\midrule
Synthesis $=$ pseudo-inverse & proved & Lemma~\ref{lem:pinv}\\
Truth is an exact fixed point & proved $+$ verified & Prop.~\ref{prop:fixed}\\
Min-overlap upper bound & proved & Thm.~\ref{thm:overlap}\\
Disjoint supports $\Rightarrow$ gauge & proved & Thm.~\ref{thm:overlap}\\
Bound tight up to $1/\sqrt2$ & measured ($r{=}0.9995$) & M.R.~\ref{meas:tight}\\
Short window pins close pairs $2\times$ & measured & M.R.~\ref{meas:pair}\\
Visibility law rectangles & support argument $+$ measured & Prop.~\ref{prop:visibility}\\
$r^{*}\in(4,8]$, interior redundancy at $r{\le}4$ & measured & M.R.~\ref{meas:rstar}\\
Ladder theorem with closed-form $r^{*}(w)$ & conjectured & Conj.~\ref{conj:ladder}\\
Strength--rate law (rate $\sim$ strength$^2$) & measured & M.R.~\ref{meas:rate}\\
Seed carries global search & measured & M.R.~\ref{meas:seed}\\
Relaxed stationarity & proved & \S\ref{sec:relaxed}\\
DR/AP gap discriminates failure class & measured & \S\ref{sec:relaxed}, \S\ref{sec:domains}\\
Complexity $\Theta(PL\log N)$ per pass & proved $+$ timed & Prop.~\ref{prop:complexity}\\
\bottomrule
\end{tabular}
\end{center}

Open items for full formality: a proof of Conjecture~\ref{conj:ladder}
from Theorem~\ref{thm:overlap} with explicit Hann support constants; a
transversality-angle theorem identifying the local convergence rate with
the min-overlap functional (upgrading Measured Result~\ref{meas:rate});
and conditions on the seed error under which one iteration attains a
prescribed fidelity (quantifying Measured Result~\ref{meas:seed}).

\begin{thebibliography}{9}
\bibitem{valdivia2026} D.~Valdivia, E.~Cazelles, C.~F\'evotte,
``Enhancing time-frequency resolution with optimal transport and
barycentric fusion of multiple spectrograms,'' arXiv:2604.15055, 2026.
\bibitem{griffinlim1984} D.~Griffin, J.~Lim, ``Signal estimation from
modified short-time Fourier transform,'' IEEE TASSP 32(2), 1984.
\bibitem{candes2013} E.~Cand\`es, T.~Strohmer, V.~Voroninski,
``PhaseLift: exact and stable signal recovery from magnitude measurements
via convex programming,'' Comm.\ Pure Appl.\ Math.\ 66(8), 2013.
\bibitem{lewisluke2009} A.~Lewis, D.~R.~Luke, J.~Malick, ``Local linear
convergence for alternating and averaged nonconvex projections,''
Found.\ Comput.\ Math.\ 9, 2009.
\bibitem{drusvyatskiy2015} D.~Drusvyatskiy, A.~Ioffe, A.~Lewis,
``Transversality and alternating projections for nonconvex sets,''
Found.\ Comput.\ Math.\ 15, 2015.
\bibitem{prusa2017} Z.~Pr\r{u}\v{s}a, P.~Balazs, P.~S{\o}ndergaard,
``A noniterative method for reconstruction of phase from STFT magnitude,''
IEEE/ACM TASLP 25(5), 2017.
\bibitem{fienup1982} J.~R.~Fienup, ``Phase retrieval algorithms: a
comparison,'' Appl.\ Opt.\ 21(15), 1982.
\bibitem{luke2005} D.~R.~Luke, ``Relaxed averaged alternating reflections
for diffraction imaging,'' Inverse Problems 21, 2005.
\end{thebibliography}

\end{document}
